\subsubsection*{2.1.3.}
现在证明 2.1.2. 数乘给出 $X$ 的恒等映射与从 $X$ 到 $X$ 取值为 $O$ 的常值映射之间的一个同伦; 这就证明了 (i).

利用乘数趋于无穷大的数乘, 可见 $Y$ 为 $X_1^*$ 的形变收缩核. 然后由五引理及下图得到断言 (ii):
\begin{center}
\begin{tikzcd}[row sep=large,column sep=large]
{} \arrow[r,"\partial"]
  & H^i_{\{O\}}(X) \arrow[r] \arrow[d]
  & H^i(X_1) \arrow[r] \arrow[d,equal]
  & H^i(X_1^*) \arrow[r,"\partial"] \arrow[d,"\wr"']
  & {} \\
{} \arrow[r,"\partial"]
  & H^i_c(X) \arrow[r]
  & H^i(X_1) \arrow[r]
  & H^i(Y) \arrow[r,"\partial"]
  & {} .
\end{tikzcd}
\end{center}

\subsubsection*{推论 2.1.4.}
\[
H^i(X^*)\xrightarrow{\ \sim\ }H^i(X^*_{(O)}).
\]
事实上, 对下图应用五引理:
\begin{center}
\begin{tikzcd}[row sep=large,column sep=large]
{} \arrow[r]
  & H^i_{\{O\}}(X) \arrow[r] \arrow[d,"\wr"]
  & H^i(X) \arrow[r]
      \arrow[d,"\wr"',"{(2.1.2(i))}"]
  & H^i(X^*) \arrow[r] \arrow[d]
  & {} \\
{} \arrow[r]
  & H^i_{\{O\}}(X_{(O)}) \arrow[r]
  & H^i(X_{(O)}) \arrow[r]
  & H^i(X^*_{(O)}) \arrow[r]
  & {} .
\end{tikzcd}
\end{center}

\subsubsection*{2.1.5.}
$h$ 在 $\widetilde X_1-Y_0$ 和 $\widetilde X_1-Y$ 上的限制的 Leray 谱序列表明, 限制映射
\[
H^i(\widetilde X_1-Y_0)\xrightarrow{\ \sim\ }H^i(Y),
\qquad
H^i(\widetilde X)\xrightarrow{\ \sim\ }H^i(Y_0)=H^i(Y)
\]
是同构. 在这些同构下, 对 $(\widetilde X_1-Y_0,Y)$ 及 $(\widetilde X_1-Y,Y_0)$ 的上同调长正合列构成下图的两行:
\begin{equation}\tag{2.1.5.1}
\begin{tikzcd}[column sep=large]
H^{i-1}(X^*) \arrow[r] \arrow[d,equal]
  & H^{i-2}(Y)(-1) \arrow[r,"\eta"'] \arrow[d,equal]
  & H^i(Y) \arrow[r] \arrow[d,dash,"-1"]
  & H^i(X^*) \arrow[r] \arrow[d,equal]
  & {} \\
H^{i-1}(X^*) \arrow[r]
  & H^{i-2}(Y)(-1) \arrow[r,"-\eta"']
  & H^i(Y) \arrow[r]
  & H^i(X^*) \arrow[r]
  & {} .
\end{tikzcd}
\end{equation}

\subsubsection*{引理 2.1.6.}
\underline{(2.1.5.1) 中以 $\eta$ 和 $-\eta$ 表示的箭头是与一个超平面截面的上同调类的杯积. 图 (2.1.5.1) 交换.}

第一个断言成立, 因为 $\eta$(分别为 $-\eta$)是 $\widetilde X_1-Y_0$(分别为 $\widetilde X_1-Y$)上可逆层 $\mathcal O(Y)$(分别为 $\mathcal O(Y_0)$)的上同调类在 $Y$(分别为 $Y_0\simeq Y$)上的限制.

使用 $h$ 的 Leray 谱序列. (2.1.5.1) 的两行是将函子 $H(Y,-)$ 应用于区分三角
\begin{center}
\begin{tikzcd}[column sep=large]
{} & y_{\infty*}Ry_\infty^!F \arrow[r]
  & R(h|\widetilde X_1-Y_0)_*F
      \arrow[r,"\textcircled{\scriptsize 1}"]
  & R(h|X^*)_*F \arrow[r,"\textcircled{\scriptsize 2}"]
  & {} \\
\text{及}
  & y_{0*}Ry_0^!F \arrow[r]
  & R(h|\widetilde X_1-Y)_*F \arrow[r]
  & R(h|X^*)_*F \arrow[r]
  & {} .
\end{tikzcd}
\end{center}
得到的. 这些三角顶点处的复形的非零上同调层分别出现在次数 $(2)$、$(0)$ 以及 $(0$ 和 $1)$ 上. 因此, 箭头 \textcircled{\scriptsize 1} 和 \textcircled{\scriptsize 2} 由其在上同调层上的作用确定. 这就将问题归结为 $Y$ 为一个点的情形, 留给读者.